% 第5章 纯仿真与域随机化的备份避障策略

\chapter{纯仿真与域随机化的备份避障策略}\label{ch:backup_policy}

\section{问题定义}\label{sec:backup_problem}

% HIL训练过程中，人手进入工作空间提供演示
% → 需要主动避碰保障安全
% → 任务策略冻结，备份策略接管
% → 人手离开后归还控制权给任务策略
%
% 需求：
% - 避免与移动人手碰撞
% - 最小化位移（原地闪避）
% - 保持方块夹持（不掉落）
% - 遵守工作空间约束（C区禁入）

\section{场景设计}\label{sec:scenarios}

% S1：单移动障碍物（28D × 3帧 = 84D）
%   - 1个移动人手，4种运动模式
%
% S2：移动+静止障碍物（38D × 3帧 = 114D）
%   - 1个移动人手 + 1个静止障碍物
%
% 每个障碍物观测（10D）：
%   [active, pos(3), vel(3), rel_pos_to_tcp(3)]
%
% 帧堆叠（3帧）：提供速度估计和运动预测的时序信息

\section{障碍物运动模式}\label{sec:motion_modes}

% 4种模式模拟不同的人手接近方式：
%
% LINEAR：匀速直线朝TCP
%   - 最简单，基线威胁模型
%
% ARC：弧线接近，每步旋转方向向量
%   - 模拟绕过障碍物的伸手动作
%
% STOP\_GO：朝TCP但随机停顿（30\%停顿概率）
%   - 模拟犹豫的人手动作
%
% PASSING：垂直于TCP方向，速度更快（1.3倍）
%   - 模拟人手快速穿过工作空间而非针对机器人
%
% 每个episode随机选择模式 → 策略必须处理所有模式

\section{奖励函数设计演化}\label{sec:backup_reward}

% 从V1到V3的演化——每次失败教会一个设计原则
%
% V1：接近度奖励（远离人手）
%   $r = -\text{collision} + \text{distance\_bonus}$
%   失败：策略学会逃跑，违背"原地闪避"目标
%   教训：奖励必须与部署目标一致
%
% V2：全负奖励
%   $r = -0.5 \|\text{displacement}\| - 0.01 \|\vect{a}\| - 0.01 \|\Delta\vect{a}\|$
%   失败："快速死亡优于缓慢死亡"
%     10步累积负奖励 > 立即碰撞惩罚
%   教训：存活必须始终优于终止
%
% V3：非负即时奖励 + 大终止惩罚（Kiemel et al. 2024）
%   每步（始终 $\geq 0$）：
%     $r = 0.5 - 0.5\|\text{displacement}\| - 0.01\|\vect{a}\| - 0.01\|\Delta\vect{a}\|$
%   终止（碰撞/掉块/C区/位移$>$15cm）：
%     $r = -10.0$
%   完成奖励：$+5.0$
%   折扣因子：$\gamma = 1.0$
%
%   期望回报：
%     完美静止：+10.0
%     小幅闪避（5cm）：+9.7
%     大幅闪避（12cm）：+8.5
%     碰撞：-10.0
%
%   设计原则：
%   1. 即时奖励必须非负
%   2. 终止惩罚 $\gg$ 累积存活奖励
%   3. 位移需要硬约束而非软惩罚
%   4. 短episode需要 $\gamma=1.0$

\section{面向Sim-to-Real的域随机化}\label{sec:dr}

% 仅作用于观测层，不影响碰撞检测
%
% | 参数         | 分布              | 真实来源               |
% |-------------|-------------------|----------------------|
% | 位置噪声     | $\mathcal{N}(0, 0.03)$  | MediaPipe + D455精度 |
% | 速度噪声     | $\mathcal{N}(0, 0.01)$  | 帧间差分抖动         |
% | 观测延迟     | $U(0, 2)$ 步            | 检测延迟             |
% | 障碍物速度   | $U(0.005, 0.015)$       | 不同手速             |
% | TCP噪声      | $\mathcal{N}(0, 0.005)$ | 外部定位噪声         |
%
% 安全原则：碰撞检测使用真实位置
% → DR使策略对噪声感知鲁棒
% → 但安全保障不受DR影响

\section{Sim-to-Real对齐分析}\label{sec:sim2real_alignment}

% 差距分析：仿真 vs 真机
% - 人手位置：MuJoCo真值 vs MediaPipe+深度 → DR覆盖
% - 人手速度：差分 vs 帧间差分 → DR覆盖
% - 观测延迟：0 vs ~10ms → DR覆盖
% - 机器人状态：无差距（编码器）
% - 控制器：OSC vs 阻抗控制 → 等价（见第7章）
% - 人手运动：4种模式 vs 真实人手 → 模式覆盖典型模式
